\documentclass[../../../include/open-logic-section]{subfiles}

\begin{document}
\olfileid[pt]{sth}{card-arithmetic}{simp}

\olsection{Simplificando Adição e Multiplicação}

Acontece que a adição e a multiplicação cardinais transfinitas são
\emph{extremamente} fáceis. Isto segue do fato de os cardinais serem
(certos) ordinais, logo bem ordenados, e poderem ser manipulados de
certa forma. Mostrar isto, porém, \emph{não} é tão fácil. Para começar,
precisamos de uma definição hábil:

\begin{defn}
Definimos uma \emph{ordenação canónica}, $\canonord$, em pares de
ordinais, estipulando que $\tuple{\alpha_1, \alpha_2} \canonord
\tuple{\beta_1, \beta_2}$ se e só se se verifica uma das seguintes condições:
\begin{enumerate}
	\item $\max(\alpha_1, \alpha_2) < \max(\beta_1, \beta_2)$; ou
	\item $\max(\alpha_1, \alpha_2) = \max(\beta_1, \beta_2)$ e
	$\alpha_1 < \beta_1$; ou
	\item $\max(\alpha_1, \alpha_2) = \max(\beta_1, \beta_2)$ e
	$\alpha_1 = \beta_1$ e $\alpha_2 < \beta_2$
\end{enumerate}
\end{defn}

\begin{lem}
$\tuple{\alpha \times \alpha, \canonord}$ é uma boa ordem, para qualquer
ordinal $\alpha$.
\end{lem}

\begin{proof}
Evidentemente $\canonord$ é conexa em $\alpha \times \alpha$. De fato,
suponhamos que nem $\tuple{\alpha_1, \alpha_2}$ nem $\tuple{\beta_1,
\beta_2}$ é menor do que a outra no sentido de $\canonord$. Então $\max(\alpha_1,
\alpha_2) = \max(\beta_1, \beta_2)$ e $\alpha_1 = \beta_1$ e
$\alpha_2 = \beta_2$, pelo que $\tuple{\alpha_1, \alpha_2} =
\tuple{\beta_1,  \beta_2}$.

Para mostrar que é boa ordem, seja $X \subseteq \alpha\times\alpha$ não
vazio. Como $\alpha$ é um ordinal, algum $\delta$ é o menor
membro de $\Setabs{\max(\gamma_1, \gamma_2)}{\tuple{\gamma_1,
\gamma_2} \in X}$. Descartemos agora todos os pares de
$\Setabs{\tuple{\gamma_1,\gamma_2} \in X}{\max(\gamma_1, \gamma_2) =
\delta}$ exceto os de primeira coordenada mínima; entre estes,
o par de segunda coordenada mínima é o elemento $\canonord$-mínimo
de $X$.
\end{proof}
\noindent
Segue-se uma observação mínima, mas simples:

\begin{prop}\ollabel{simplecardproduct}
Se $\cardeq{\alpha}{\beta}$, então $\cardeq{\alpha \times \alpha}{\beta
\times \beta}$. 
\end{prop}

\begin{proof}
Basta deixar que $f \colon \alpha \to \beta$ induza $\tuple{\gamma_1,
\gamma_2} \mapsto \tuple{f(\gamma_1), f(\gamma_2)}$.
\end{proof}
\noindent
E agora aplicaremos tudo isto para provar um lema crucial:
\begin{lem}\ollabel{alphatimesalpha}
$\cardeq{\alpha}{\alpha \times \alpha}$, para qualquer ordinal infinito
$\alpha$
\end{lem}

\begin{proof}
Por redução ao absurdo, seja $\alpha$ o menor ordinal infinito para o qual
isto falha. \olref[sfr][siz][zigzag]{natsquaredenumerable} mostra
que $\cardeq{\omega}{\omega\times\omega}$, logo $\omega \in \alpha$.
Além disso, $\alpha$ é um cardinal: suponhamos o contrário, por redução ao absurdo;
então $\card{\alpha} \in \alpha$, pelo que
$\cardeq{\card{\alpha}}{\card{\alpha} \times \card{\alpha}}$, por
hipótese; e $\cardeq{\card{\alpha}}{\alpha}$ por definição; logo
$\cardeq{\alpha}{\alpha\times\alpha}$ por
\olref{simplecardproduct}. 

Ora, para cada $\tuple{\gamma_1, \gamma_2} \in \alpha \times \alpha$,
consideremos o segmento:
\begin{align*}
	\text{Seg}(\gamma_1, \gamma_2) &= \Setabs{\tuple{\delta_1, \delta_2} \in \alpha \times \alpha}{\tuple{\delta_1, \delta_2} \canonord \tuple{\gamma_1, \gamma_2}}
\end{align*}
Seja $\gamma = \max(\gamma_1, \gamma_2)$; note-se que $\tuple{\gamma_1, \gamma_2} \canonord \tuple{\gamma+1, \gamma + 1}$. Assim, quando $\gamma$ é infinito, observe-se:
\begin{align*}
	\text{Seg}(\gamma_1, \gamma_2) & 
	\precsim ((\gamma \ordplus 1)\ordtimes (\gamma \ordplus 1))\\
	&\approx (\gamma \ordtimes \gamma)
	\text{, por \olref[ord-arithmetic][using-addition]{ordinfinitycharacter} e 
	\olref{simplecardproduct}}\\
	&\approx \gamma \text{, pela hipótese de indução}\\
	& \prec \alpha\text{, pois $\alpha$ é um cardinal}
\end{align*}
Logo $\ordtype{\alpha\times \alpha, \canonord} \leq \alpha$, e portanto
$\cardle{\alpha \times \alpha}{\alpha}$. Como obviamente
$\cardle{\alpha}{\alpha \times \alpha}$, o resultado segue por
Schr\"oder-Bernstein. 
\end{proof}

Finalmente, chegamos ao nosso resultado simplificador:

\begin{thm}\ollabel{cardplustimesmax}
Se $\cardfont{a}, \cardfont{b}$ são cardinais infinitos, então:
\[
	\cardfont{a}
\cardtimes \cardfont{b} = \cardfont{a} \cardplus \cardfont{b} =
\text{max}(\cardfont{a}, \cardfont{b}).
\]
\end{thm}

\begin{proof}
Sem perda de generalidade, suponhamos $\cardfont{a} = \max(\cardfont{a},
\cardfont{b})$. Então, invocando \olref{alphatimesalpha},
$\cardfont{a}\cardtimes\cardfont{a} = \cardfont{a} \leq \cardfont{a}
\cardplus \cardfont{b} \leq \cardfont{a} \cardplus \cardfont{a} \leq
\cardfont{a} \cardtimes \cardfont{a}$. \end{proof}\noindent Do mesmo modo,
se $\cardfont{a}$ é infinito, uma união de cardinalidade $\cardfont{a}$ de
conjuntos de cardinalidade $\leq\cardfont{a}$ tem cardinalidade $\leq\cardfont{a}$:

\begin{prop}\ollabel{kappaunionkappasize}
Seja $\cardfont{a}$ um cardinal infinito. Para cada ordinal $\beta
\in \cardfont{a}$, seja $X_\beta$ um conjunto com $\card{X_\beta} \leq
\cardfont{a}$. Então $\card{\bigcup_{\beta \in \cardfont{a}} X_\beta}
\leq \cardfont{a}$.
\end{prop}

\begin{proof}
Para cada $\beta \in \cardfont{a}$, fixe !!a{injection} $f_\beta \colon
X_\beta \to \cardfont{a}$.\footnote{Como são estes ``fixos''? Ver \olref[sth][choice][countablechoice]{sec}.} Defina !!a{injection} $g \colon
\bigcup_{\beta \in \cardfont{a}} X_\beta \to \cardfont{a} \times
\cardfont{a}$ por $g(v) = \tuple{\beta, f_\beta(v)}$, onde $v \in
X_\beta$ e $v \notin X_\gamma$ para qualquer $\gamma \in \beta$. Ora,
$\bigcup_{\beta \in \cardfont{a}} X_\beta \preceq \cardfont{a} \times
\cardfont{a} \approx \cardfont{a}$ por \olref{cardplustimesmax}.
\end{proof}

\end{document}
